% Abstract
Current RLHF optimizes language models for helpfulness alone, neglecting the Human$\to$AI direction of alignment: controllability. We introduce \textbf{bidirectional alignment}, augmenting standard helpfulness rewards with IFEval-derived controllability signals via the combined reward $R = \alpha \cdot R_{\text{helpfulness}} + \beta \cdot R_{\text{IFEval}}$. We convert discrete constraint checks into differentiable rewards using sigmoid soft thresholds.

Experiments on Llama-3-8B-Instruct validate three predictions: (1) bidirectional models achieve \textbf{+3.4pp} held-out IFEval strict accuracy over helpfulness-only baselines; (2) at $\alpha=0.8$, models retain \textbf{96.4\%} of baseline AlpacaEval performance; (3) explicit constraint training transfers to implicit safety constraints, with \textbf{+2.7pp} TruthfulQA and \textbf{+4.2pp} BBQ improvements and strong positive correlation (\textbf{r=0.94}) between IFEval gains and safety gains.

Our most surprising finding is the explicit$\to$implicit transfer: models trained on format and length constraints improve on safety benchmarks measuring truthfulness and bias---constraints never seen during training. We estimate a $\sim$15\% transfer coefficient: each 1pp IFEval improvement predicts $\sim$0.15pp safety improvement.

This work provides the first empirical validation of the bidirectional alignment framework (Sun et al. 2024), introduces IFEval as a differentiable training signal, and characterizes the helpfulness-controllability Pareto frontier.
